\section{Local convergence: geometry that survives near a minimizer}\label{art:local}\label{art:positive}

The global theorem requires the initial sublevel to lie in the region where
its representation is valid. A local theorem must choose that region before
seeing a BFGS run or its initial metric. Convexity and a unique minimizer
supply precisely this localization. We first use it for finite polyhedral
pieces, where inward accuracy has a direct geometric explanation, and then
for curved planar switching interfaces, where a moving comparison direction
is essential to the proof.

\subsection{Localization and finite polyhedral pieces}

The distinction between a local objective hypothesis and a hypothesis on the
whole initial sublevel is handled by a simple convex trapping argument.

\begin{lemma}[A common local basin]\label{art:localization}
Let $F:\R^n\to\R$ be finite and convex with a unique minimizer $x_*$.
Every sublevel is compact. For each neighborhood $U$ of $x_*$, there is
$\varepsilon>0$ such that $\{F\le F(x_*)+\varepsilon\}\subset U$.
Thus any admissible run starting in that sublevel stays in $U$, independently
of its initial SPD matrix.
\end{lemma}
\begin{proof}
Choose $r>0$ with $\overline B_r(x_*)\subset U$.
Continuity, compactness of the sphere and uniqueness give
\[
 \omega=\min_{\norm{u}=r}\bigl(F(x_*+u)-F(x_*)\bigr)>0.
\]
For $\norm{x-x_*}\ge r$, convexity at the point at distance $r$ on the
segment from $x_*$ to $x$ gives
\[
 F(x)-F(x_*)\ge\frac{\norm{x-x_*}}r\omega.
\]
Consequently $F$ is coercive, hence has compact sublevels, and
$\varepsilon=\omega/2$ gives the asserted inclusion. Armijo preserves it.
\end{proof}

A finite local partition controls which formulas an inward comparison segment
can cross. After removing the finitely many cells that miss $x_*$, every
remaining cell contains $x_*$. Convexity of a cell then keeps the entire
inward segment in that same cell. Its smooth extension yields quadratic
model accuracy even if the full gradient jumps across cell boundaries.
The algorithm need not identify a cell or a manifold; only the comparison
segment uses the cell geometry.

\begin{theorem}[Local convergence on finite polyhedral cells]\label{art:polyhedral}
Let $n\ge1$ and let $F:\R^n\to\R$ be finite and convex, with a unique minimizer $x_*$.
Suppose an open neighborhood $U$ of $x_*$ is covered by finitely many closed
full-dimensional polyhedra $P_i$. For each $i$, suppose there is an open set
$V_i\supset U\cap P_i$ and a function $f_i\in C^{1,1}_{\rm loc}(V_i)$ such that
\[
 F(x)=f_i(x)\qquad(x\in U\cap P_i).
\]
Then there exists $\varepsilon>0$ such that every infinite full-BFGS run with
arbitrary $H_0\succ0$, differentiable accepted iterates, fixed
Armijo--weak-Wolfe constants $0<c_1<c_2<1$, and
$F(x_0)\le F(x_*)+\varepsilon$ satisfies
\[
 x_k\in U\quad(k\ge0),\qquad F(x_k)\to F(x_*),\qquad x_k\to x_*.
\]
The basin is chosen before the run and its initial matrix. Only the full
objective must be convex; the extensions need not be convex outside their cells.
\end{theorem}


\begin{proof}[Theorem~\ref{art:polyhedral}]
First discard cells not containing $x_*$. Each such closed cell has an open
complement containing $x_*$; finiteness supplies one ball that avoids them all.
For every remaining cell, its extension domain $V_i$ contains $x_*$.
Shrink to a common ball $B_r(x_*)\subset U\cap\bigcap_{i:x_*\in P_i}V_i$,
and then shrink again so that the gradients of all remaining extensions have
one Lipschitz bound $L$ there. This uses only finitely many local extensions;
they need not have shared a domain before shrinking.

Let $x$ be a differentiability point of $F$ in this ball, and choose a cell
$P_i$ containing it. Convexity of $P_i$ places the whole inward segment
$x-\tau(x-x_*)$, $0\le\tau\le1$, in $P_i$. Both $F$ and $f_i$ agree
along that segment. Comparing their right derivatives at $\tau=0$ gives
\[
 \nabla F(x)^T(x-x_*)=\nabla f_i(x)^T(x-x_*).
\]
Only equality in the radial direction is needed, including at cell boundaries.
The smooth Taylor estimate along the same segment gives
\[
 0\le F(x-\tau r_x)-F(x)+\tau\nabla F(x)^Tr_x
       \le\tfrac12L\tau^2\norm{r_x}^2
       \le\tfrac12Lr^2\tau^2.
\]
The first inequality uses convexity of $F$, not convexity of the extension.
Lemma~\ref{art:localization} supplies a compact positive sublevel wholly
contained in this ball, chosen before the run or its matrix. The inward
condition holds on that sublevel, so Corollary~\ref{art:radial-convergence}
gives value and iterate convergence. Armijo gives the asserted containment.
\end{proof}

\begin{corollary}[Smooth losses with polyhedral or quadratic pieces]\label{art:plq}
A finite convex piecewise linear-quadratic objective with a unique minimizer
has the local convergence property of Theorem~\ref{art:polyhedral}.
The same holds for a convex locally $C^{1,1}$ loss plus a finite convex
piecewise linear-quadratic regularizer, provided the sum has a unique minimizer.
\end{corollary}
\begin{proof}
By piecewise linear-quadratic we mean a representation by quadratic polynomials
on a finite cover by closed full-dimensional polyhedral cells. Each polynomial is a smooth
extension of its cell formula; adding the smooth loss preserves the required
local $C^{1,1}$ extensions. Theorem~\ref{art:polyhedral} applies.
\end{proof}

In particular, a finite convex polyhedral penalty, such as a finite sum of
absolute-value or hinge terms, is piecewise affine on a finite polyhedral
partition. It therefore falls under this corollary when added to a locally
$C^{1,1}$ convex loss,
under the stated uniqueness assumption. Different cells may have different
quadratic Hessians. A finite maximum of arbitrary quadratic formulas need not
admit such a polyhedral partition and is not automatically covered.
There is also no claim that iterates identify a cell or that rejected trials
remain in the local structural neighborhood.


The finiteness and cell geometry are substantive. Finitely many curved
switching interfaces need not preserve inward segments, and the negative
examples later in the paper show that regularity of smooth branches alone
cannot replace the stated geometric hypothesis. The theorem gives a
sufficient local basin, without claiming necessity or a uniform rate.

\begin{corollary}[A local centered-composite basin]\label{art:local-centered}
If the representation and smoothness assumptions of
Theorem~\ref{art:centered} hold on a neighborhood of a unique minimizer $x_*$,
then some $\varepsilon>0$, chosen from the objective alone, has the following
property: every admissible run starting with
$F(x_0)\le F(x_*)+\varepsilon$ stays in that neighborhood and converges to
$x_*$ in value and iterates.
\end{corollary}
\begin{proof}
Shrink to an open convex neighborhood where the representation holds.
Lemma~\ref{art:localization} gives a compact sublevel contained there.
Armijo traps the accepted points, and Theorem~\ref{art:centered} applies.
\end{proof}

For example, a locally $C^{1,1}$ convex loss plus finitely many Euclidean
residual norms $\norm{A_jx-b_j}$ has this local form. A residual nonzero at
$x_*$ remains nonzero nearby, so its Euclidean norm is smooth there; a zero
residual gives the centered sublinear function $\norm{A_j(x-x_*)}$.
Summing these decompositions proves the representation. Under uniqueness,
the corollary applies. This example allows nonpolyhedral penalties while
retaining the same radial mechanism.

\subsection{A moving direction for curved planar switching}

The radial direction can fail for a different reason: a curved interface may
carry a gradient jump with a nonzero radial component. The secant correction
then blows up as endpoints approach opposite sides of the interface.
The transported framework allows the comparison direction to follow the
interface instead. It must first remain visible to every local supporting
model, then cancel the leading jump, and finally move at a cost controlled
by the actual objective decrease. The field is constructed from the
objective; BFGS never evaluates it.

For $F:\R^2\to\R$, write
\[
 \mathcal C_{\rm aff}=
 \{p:0\in\aff\partial F(p)\}.
\]
At a point outside this set, let $a$ be the least-norm point of
$\aff\partial F(p)$. Orthogonality to its direction space gives
$g^Ta=\norm{a}^2$ for every $g\in\partial F(p)$. Thus
$-a/\norm{a}^2$ decreases all supporting models at the same unit rate.
Affine criticality concerns failure of this common-rate direction;
ordinary stationarity concerns failure of every strict common descent
direction. They need not coincide.

The two planar hypotheses have separate roles. Strictly convex positive
sublevels prohibit long accepted chords with vanishing value decrease.
Regular two-branch switching makes each positive-level corner move smoothly
with the level and supplies a direction paired equally with both gradients.
Neither condition imposes a matrix limit or a lower bound on all eigenvalues
of a Hessian. A common positive lower Hessian bound is one sufficient way
to obtain strict contours, but is not required.

\begin{theorem}[An affine-critical cluster point in the plane]\label{art:planar-cluster}
Let $F:\R^2\to\R$ be finite and convex, attain its minimum $F_*$, and have
a compact sublevel $K=\{F\le b\}$ with $b>F_*$. Assume the following
objective geometry:

\begin{enumerate}
\item For every $F_*<t\le b$, the boundary of $\{F\le t\}$ contains no
nontrivial line segment.
\item At every $p\in K\setminus\mathcal C_{\rm aff}$, either $F$ is $C^2$
in a neighborhood of $p$, or it equals $\max\{f_1,f_2\}$ there, where
$f_1,f_2$ are convex $C^2$ functions, both active at $p$, and
$\nabla f_1(p),\nabla f_2(p)$ are linearly independent.
\end{enumerate}
Every infinite admissible full-BFGS run starting in $K$ has an affine-critical
cluster point. No local representation is required at affine-critical points.
\end{theorem}
\begin{proof}
The values decrease to $\ell\ge F_*$, the gradients are bounded on $K$, and
$\sum d_k<\infty$. Let $\Gamma$ be the nonempty compact cluster set.
Every point of $\Gamma$ has value $\ell$. If $\ell=F_*$, its points are
minimizers and hence affine critical. Suppose, for a contradiction, that
$\Gamma$ contains no affine-critical point. Then $\ell>F_*$.

The steps tend to zero. Otherwise compactness gives a subsequence with
$x_k\to x$, $x_{k+1}\to z\ne x$. For $0\le\theta\le1$, the supporting
inequality and convexity give
\[
 F(x_k)-\theta d_k
 \le F((1-\theta)x_k+\theta x_{k+1})
 \le(1-\theta)F(x_k)+\theta F(x_{k+1}).
\]
Passing to the limit shows $F=\ell$ on $[x,z]$. Such level points lie on
the boundary of $\{F\le\ell\}$: if one were interior, extending the line
from a minimizer slightly beyond that point and applying convexity would
give value strictly less than $\ell$. The resulting boundary segment
contradicts the first hypothesis.

Lemma~\ref{art:planar-field}, proved in Appendix~\ref{art:planar-appendix},
constructs a continuous comparison field $w$ on a neighborhood of the whole
set $\Gamma$. For each sufficiently close eligible pair in that neighborhood,
it gives uniform constants $m>0,C,L\ge0$ and scalars $\lambda>0,\beta$ with
\[
 \begin{gathered}
 (g^Tw(x))^2\ge m,\qquad (y^Tw(z))^2/t\le C,\\
 w(z)=\lambda w(x)+\beta s,\qquad \lambda^2\le1+Ld.
 \end{gathered}
\]
Compactness implies $\operatorname{dist}(x_k,\Gamma)\to0$, so these estimates
apply to the entire tail, not just a selected subsequence.
With $E_k=w(x_k)^TB_kw(x_k)$, the Hessian update and $M_ks_k=0$ yield
\[
 E_{k+1}=
 \lambda_k^2\left(E_k-\frac{(g_k^Tw(x_k))^2}{q_k}\right)
       +\frac{(y_k^Tw(x_{k+1}))^2}{t_k}
 \le(1+Ld_k)(E_k-m/q_k)+C.
\]
The parenthesized quadratic form is nonnegative. Dividing by
$\prod_{j=J}^{k}(1+Ld_j)$ and summing, exactly as in the proof of
Theorem~\ref{art:certificate}, bounds $\sum_{k=J}^{N-1}q_k^{-1}$ by a constant
times $N-J+1$. This contradicts Theorem~\ref{art:reciprocal-growth}.
\end{proof}

\begin{corollary}[When affine criticality is correct]\label{art:planar-convergence}
Under the hypotheses of Theorem~\ref{art:planar-cluster}, assume in addition
that $\mathcal C_{\rm aff}\cap K\subseteq\arg\min F$. Every admissible run
starting in $K$ then converges in value to $F_*$ and in distance to
$\arg\min F$. A unique minimizer gives convergence of the entire sequence.
\end{corollary}
\begin{proof}
The theorem supplies an affine-critical cluster point, which is optimal by
hypothesis. Its value equals the limiting value of the decreasing objective
sequence. Hence that value is $F_*$, and compactness gives the distance and
uniqueness conclusions as in Theorem~\ref{art:certificate}.
\end{proof}

This condition says that nonoptimal points retain equal-rate descent, which
is stronger than ordinary descent. It is a sufficient geometric route in the
stated planar setting, not a universal criterion for either convergence or
failure. The strict-contour and regular-switching assumptions are independent
parts of this route. In particular, the argument neither handles arbitrary
flat contours nor transfers automatically to dimension three: the tangent
line used in the transport construction is then a higher-dimensional space.

\subsection{Nondegenerate smooth maxima and an explicit nonradial basin}

The extra affine-correctness condition can be verified from the active
models at the minimizer. Positive optimal multipliers ensure that an affine
balance cannot lose a branch nearby without losing stationarity; affine
independence makes that balance unique and continuous. In the plane these
same assumptions leave at most two active branches at a noncritical point,
which supplies the required regular switching.

\begin{corollary}[Nondegenerate planar smooth maxima]\label{art:planar-max}
Let $F:\R^2\to\R$ be finite and convex with a unique minimizer $x_*$. Near
$x_*$ suppose $F=\max_{1\le i\le m}f_i$, where the branches are convex and
$C^2$. Suppose sufficiently low positive sublevels are strictly convex in
the boundary-segment sense of Theorem~\ref{art:planar-cluster}. For the
active set $I_*$ at $x_*$, suppose the gradients are affinely independent
and admit weights
\[
 \lambda_i^*>0,\qquad \sum_{i\in I_*}\lambda_i^*=1,
 \qquad \sum_{i\in I_*}\lambda_i^*\nabla f_i(x_*)=0.
\]
Then some $\varepsilon>0$, chosen independently of $H_0$ and the valid
line-search constants, satisfies
\[
 F(x_0)<F_*+\varepsilon
 \quad\Longrightarrow\quad F(x_k)\to F_*,\qquad x_k\to x_*
\]
for every infinite admissible full-BFGS run. A common positive lower Hessian
bound on the locally relevant branches suffices for the contour hypothesis.
\end{corollary}
\begin{proof}
Discard the branches inactive at $x_*$ by their strict value gaps. Let $A(x)$
have columns $(\nabla f_i(x),1)$ for $i\in I_*$. Affine independence makes
$A(x_*)$ full column rank, and this persists nearby. If
$A(p)\lambda=(0,0,1)^T$ is solvable, its unique solution is
\[
 \lambda(p)=(A(p)^TA(p))^{-1}A(p)^T(0,0,1)^T.
\]
This expression is continuous and equals $\lambda^*$ at $x_*$. Shrink the
neighborhood until all its coordinates are positive. The subdifferential
of a finite smooth maximum is the convex hull of its active gradients.
Indeed, differentiating the finite maximum in a direction gives the maximum
of the active gradient pairings. Every convex combination of those gradients
supports the objective locally; convexity extends the supporting inequality
globally by taking a short initial segment toward any proposed violating
point. Conversely, a subgradient outside their convex
hull could be strictly separated from that hull by a direction, contradicting
the directional-derivative inequality for a subgradient.
An affine representation of zero using gradients active at $p$, extended by
zero on inactive indices, must equal $\lambda(p)$. It therefore uses every
index and is a convex representation of zero. Thus every nearby
affine-critical point is a minimizer.

At most three planar gradients are affinely independent. Away from affine
criticality, three cannot be active, since their affine hull would be all of
$\R^2$. One active branch gives a smooth chart. For two active branches their
gradients are distinct; if they were linearly dependent, their affine line
would contain zero. Thus the gradients are independent, and the other
branches have strict gaps. This gives the required convex two-branch chart.

Lemma~\ref{art:localization} now puts a sufficiently low compact sublevel
inside the chosen neighborhood, with its upper level also below the threshold
for strict contours. Apply Corollary~\ref{art:planar-convergence}.
If all relevant branch Hessians have a common positive lower bound on a
convex neighborhood, subtracting the corresponding positive quadratic
makes every branch convex. Their maximum is therefore strongly convex
there. A sufficiently low sublevel lies in that neighborhood and cannot have
a nontrivial boundary segment, proving the last assertion.
\end{proof}

The multiplier geometry is related to the nondegeneracy used in local bundle
Newton analysis \preprintcite[Propositions~5.5--5.6]{LW} and in Gebken's
criticality results \cite{BG}.
Here the full-BFGS conclusion is derived without assumed iterate or matrix
asymptotics. The contour and planar restrictions remain explicit.

\begin{theorem}[Convergence with curved, nonradial switching]\label{art:nonradial-example}
For
\[
 \begin{aligned}
 f_0(u,v)&=u+u^2+\tfrac12v^2,\\
 f_1(u,v)&=v+\tfrac12u^2+\tfrac32v^2,\\
 f_2(u,v)&=-u-v+\tfrac32u^2+v^2,\qquad F=\max_i f_i,
 \end{aligned}
\]
every infinite admissible full-BFGS run with $F(x_0)<1/100$ converges in
value and iterates to zero, for arbitrary $H_0\succ0$ and fixed valid
Armijo--weak-Wolfe constants. Nevertheless, the radial secant condition
\eqref{art:radial-condition} fails on a positive band in every neighborhood
of zero. The obstructing endpoint pairs can themselves be chosen eligible
for $c_1=1/4,c_2=1/2$.
\end{theorem}
\begin{proof}
The branch Hessians lie between $I$ and $3I$, and
$(f_0+f_1+f_2)/3=u^2+v^2$. Thus $F\ge u^2+v^2$, zero is the unique
minimizer, and all positive contours are strict. The active gradients at
zero are $(1,0),(0,1),(-1,-1)$, with positive weights $1/3$.
To obtain the stated basin rather than an unspecified one, put
$D=3(1+v-uv)$, the determinant of the augmented gradient matrix.
Direct solution of that system gives
\[
 \begin{aligned}
 D\lambda_0&=1-4u+3v-7uv,&
 D\lambda_1&=1+2u-3v-uv,\\
 D\lambda_2&=1+2u+3v+5uv.&&
 \end{aligned}
\]
On $|u|,|v|\le1/10$, $D\ge267/100$ and the three numerators are at least
$23/100$, $49/100$, and $45/100$. The coefficients sum to one and their
weighted gradient is zero. Thus affine independence and affine correctness
hold throughout this box, by the same unique-coefficient argument as above.
The planar regularity argument also applies throughout the box. Since
$\{F\le1/100\}$ lies there, Corollary~\ref{art:planar-convergence} proves
the exact convergence claim.

For the radial obstruction choose a sufficiently small fixed $u>0$ and put
\[
 v=v(u)=\frac{-1+\sqrt{1+4u+2u^2}}2
       =u-\tfrac12u^2+O(u^3).
\]
This is the positive active tie $f_0=f_1$, with $f_2$ inactive. At that point
\[
 (\nabla f_0-\nabla f_1)^T(u,v)
 =(1+u)u-(1+2v)v=\tfrac12u^2-v^2\ne0
\]
for small $u$, as follows from the displayed expansion. Take
$x=(u,v+e)$ and $z=(u,v-e)$ with $e>0$ tending to zero. Branches 1 and 0,
respectively, are uniquely active. A direct calculation gives
\[
 \begin{aligned}
 d&=2e(1+3v+3e),&t&=2e(1+2v+4e),\\
 \Delta&=e(1+4v)+e^2.&&
 \end{aligned}
\]
As $u,e\to0^+$, $\Delta/d\to1/2$ and $t/d\to1$, so for fixed sufficiently
small $u$ and all sufficiently small $e$, the pairs are eligible for the
stated constants. Their radial numerator $y^Tz$ tends to
$u^2/2-v^2\ne0$ while $t\to0^+$. The quotient $(y^Tz)^2/t$ is unbounded.
For each chosen $u$ both endpoint values stay above half the positive tie
value; hence one fixed positive band precedes the unbounded sequence.
Choosing $u$ arbitrarily small places that band and those pairs in any
given neighborhood of zero.
\end{proof}

The example identifies the extra freedom supplied by moving geometry.
Curvature of the interface produces a radial jump, while the contour-based
field cancels it and charges its own motion to the accepted decrease.
The proof uses the original coordinates and the unmodified BFGS update;
a nonlinear change of variables or a replacement by polyhedral cells would
not establish this conclusion.
